\documentclass{article}
\usepackage{graphicx} % Required for inserting images
\usepackage{amsmath}

\title{Exercise 4 - Theoretical Questions\\
Question 1: Lazy Lists}
\author{
IDO KOZIHI \\
IDO UZIYAHU
}
\date{June 2026}

\begin{document}

\maketitle

\section*{Question 1.1: Lazy Lists}

\subsection*{(a) Equivalence Criterion}

Two lazy lists are considered equivalent if, after any finite number of applications of \texttt{tail}, either both lists are empty, or both lists have a defined \texttt{head} and these two \texttt{head} values are equal.

\subsection*{(b) Equivalence of \texttt{even-square-1} and \texttt{even-square-2}}

Let
\[
L_1 = \texttt{even-square-1}
\]
and
\[
L_2 = \texttt{even-square-2}.
\]

We consider an arbitrary natural number \(T\) produced by \texttt{integers-from 0}.

In \(L_1\), the number \(T\) is first mapped to \(T^2\), and then the filter checks whether \(T^2\) is even.

In \(L_2\), the number \(T\) is first filtered according to whether \(T\) is even, and only then it is mapped to \(T^2\).

We split into two cases.

If \(T\) is even, then \(T^2\) is also even. Therefore, in \(L_1\), the value \(T^2\) passes the filter. In \(L_2\), the value \(T\) passes the filter and is then mapped to \(T^2\). Thus, in both lists, the same value \(T^2\) appears.

If \(T\) is odd, then \(T^2\) is also odd. Therefore, in \(L_1\), the value \(T^2\) does not pass the filter. In \(L_2\), the value \(T\) does not pass the filter before the mapping. Thus, in both lists, \(T\) contributes no value to the final lazy list.

Hence, for every natural number \(T\), either both lists include the same value \(T^2\), or both lists exclude it. Since both \texttt{map} and \texttt{filter} preserve the original order of the remaining elements, both lazy lists produce the same sequence:
\[
0, 4, 16, 36, 64, \ldots
\]

Therefore, according to the equivalence criterion defined in part (a), \(L_1\) and \(L_2\) are equivalent lazy lists.


\vspace{20pt}
\section*{Question 2: CPS}
\subsection*{(a)}
A procedure \texttt{func} of type \([A_1*...*A_n \rightarrow T \mid \textit{'fail}]\) and its Success-Fail-Continuations version \texttt{func\$} of type \([A_1*...*A_n*[T \rightarrow T1]*[Empty \rightarrow T2] \rightarrow T1 \mid T2]\) are equivalent if and only if for every valid sequence \(x_1, \dots, x_n\) and any continuation functions \texttt{success} and \texttt{fail}


\begin{enumerate}
    \item \((\texttt{func} \ x_1 \dots x_n)\) will terminate and return value \(\text{res} \neq \textit{'fail} \iff (\texttt{func\$} \ x_1 \dots x_n \text{ success } \text{fail})\) results in the evaluation of \((\text{success } \text{res})\).
    
    \item \((\texttt{func} \ x_1 \dots x_n)\) will terminate and return \(\textit{'fail} \iff (\texttt{func\$} \ x_1 \dots x_n \text{ success } \text{fail})\) results in the evaluation of \((\text{fail})\).
    
    \item \((\texttt{func} \ x_1 \dots x_n)\) does not terminate \(\iff (\texttt{func\$}  \ x_1 \dots x_n \text{ success } \text{fail})\) does not terminate.
\end{enumerate}

\subsection*{2d}
We will prove \texttt{get-value} and \texttt{get-value\$} equivalence using induction on \(n = \text{len}(\text{assoc-list})\).

\subsubsection*{Base Case}
On \(n=0\), \((\text{empty? } \text{assoc-list})\) evaluates to \texttt{true}. Therefore, \texttt{get-value} will terminate and return \(\textit{'fail}\), and \texttt{get-value\$} will evaluate \((\text{fail})\).

\subsubsection*{Proposition}
For every \(n > 0\), we will assume that for every \texttt{assoc-list} with \(\text{len}(\text{assoc-list}) = n\), \texttt{get-value} and \texttt{get-value\$} are equivalent.

\subsubsection*{Induction Step}
We will prove that \((\text{get-value } \text{assoc-list } \text{key})\) and \((\text{get-value\$} \text{ assoc-list } \text{key } \text{success } \text{fail})\) are equivalent for an \texttt{assoc-list} with \(\text{len}(\text{assoc-list}) = n+1\).

For easier reading, we will refer to \((\text{get-value } \text{assoc-list } \text{key})\) as \texttt{get-value} and \((\text{get-value\$} \text{ assoc-list } \text{key } \text{success } \text{fail})\) as \texttt{get-value\$}.

Since \(\text{len}(\text{assoc-list}) > 0\), \texttt{get-value} will evaluate to:
\begin{verbatim}
(if (eq? key (car (car assoc-list)))
    (cdr (car assoc-list))
    (get-value (cdr assoc-list) key))
\end{verbatim}
and \texttt{get-value\$} will evaluate to:
\begin{verbatim}
(if (eq? key (car (car assoc-list)))
    (success (cdr (car assoc-list)))
    (get-value$ (cdr assoc-list) key success fail))
\end{verbatim}

\paragraph{Case 1:} \texttt{(eq? key (car (car assoc-list)))} evaluates to \texttt{\#t}\\
In this case, \texttt{get-value} will evaluate to \texttt{(cdr (car assoc-list))}, and \texttt{get-value\$} will evaluate to \texttt{(success (cdr (car assoc-list)))}. Therefore, they are equivalent.

\paragraph{Case 2:} \texttt{(eq? key (car (car assoc-list)))} evaluates to \texttt{\#f}\\
Here, \texttt{get-value} will return the evaluation of \texttt{(get-value (cdr assoc-list) key)}, and \texttt{get-value\$} will evaluate \texttt{(get-value\$ (cdr assoc-list) key success fail)}.

By assuming the induction hypothesis for \texttt{(cdr assoc-list)}, which has length \(n\), it follows that \texttt{(get-value (cdr assoc-list) key)} and \texttt{(get-value\$ (cdr assoc-list) key success fail)} are equivalent. Therefore, \texttt{get-value} and \texttt{get-value\$} are equivalent.

\end{document}